\section{Bayesian Latent Factor Model for Composite Quality Scores}
\label{sec:blf}

\subsection{Motivation}

A fundamental challenge in multi-model routing systems is constructing task-specific quality scores when benchmark coverage is incomplete and heterogeneous. Consider coding quality assessment: while HumanEval \cite{chen2021codex} provides function-level generation scores, LiveCodeBench \cite{jain2024livecodebench} measures competitive programming ability, and SciCode \cite{tian2024scicode} evaluates scientific computing---yet no single model is evaluated on all benchmarks. Naive approaches fail:

\begin{itemize}
    \item \textbf{Single-benchmark selection} ignores complementary signals and reduces coverage (e.g., only 73\% of models have LiveCodeBench scores).
    \item \textbf{Arithmetic mean} treats all benchmarks equally, despite varying difficulty, scale, and measurement noise.
    \item \textbf{Weighted z-scores} improve upon arithmetic mean but cannot handle missing data and ignore measurement uncertainty.
    \item \textbf{Listwise deletion} (dropping models with incomplete data) discards 60-85\% of models depending on the task.
\end{itemize}

We introduce a \textbf{Bayesian Latent Factor (BLF)} model that: (1) learns benchmark-specific weights from data rather than requiring manual specification, (2) handles missing benchmarks via principled imputation using auxiliary covariates, (3) quantifies uncertainty in composite scores, and (4) is robust to outliers and varying scales.

\subsection{Model Specification}

Let $z_{i,b}$ denote the standardized score for model $i \in \{1, \ldots, N\}$ on benchmark $b \in \{1, \ldots, B\}$. We model observed scores as arising from a latent composite quality $\theta_i$ with benchmark-specific characteristics:

\begin{align}
z_{i,b} &\sim \mathcal{N}(\alpha_b + \lambda_b \theta_i, \sigma_b^2) \label{eq:likelihood} \\
\theta_i &\sim \mathcal{N}(0, 1) \label{eq:latent_prior} \\
\alpha_b &\sim \mathcal{N}(0, 2^2) \label{eq:intercept_prior} \\
\lambda_b &\sim \text{HalfNormal}(1) \label{eq:loading_prior} \\
\sigma_b &\sim \text{HalfNormal}(1) \label{eq:noise_prior}
\end{align}

where:
\begin{itemize}
    \item $\theta_i$ is the latent composite score for model $i$ (standardized, mean 0, sd 1)
    \item $\alpha_b$ is the benchmark-specific intercept (accounts for difficulty)
    \item $\lambda_b$ is the benchmark-specific loading (learned weight, analogous to factor analysis)
    \item $\sigma_b$ is the benchmark-specific residual noise (measurement error)
\end{itemize}

This is a Bayesian factor analysis model where the latent factor $\theta_i$ represents the ``true'' composite quality, and observed benchmarks are noisy, scaled indicators of this latent construct. We assume a unidimensional latent quality construct (positive manifold), enforced by the positive constraints on loadings. The model naturally handles missing data: models with $z_{i,b} = \text{NA}$ contribute information through other observed benchmarks and the prior.

\subsection{Auxiliary Benchmarks for Covariance-Based Imputation}

For models missing primary benchmarks, we leverage high-coverage auxiliary metrics. For example, in the Composite Coding Score (CCS), we include Intelligence Index (99\% coverage, $r=0.96$ with LiveCodeBench) as an auxiliary benchmark. The model estimates $\theta_i$ for these models by exploiting the covariance structure:

\begin{equation}
p(\theta_i | z_{i,\text{aux}}) \propto p(z_{i,\text{aux}} | \theta_i) p(\theta_i)
\end{equation}

This borrowing of statistical strength is analogous to multiple imputation \cite{rubin1987multiple}, but occurs naturally within the Bayesian framework without requiring separate imputation steps.

\subsection{Inference and Posterior Computation}

We perform full Bayesian inference using Hamiltonian Monte Carlo (HMC) via PyMC \cite{abril2023pymc}:

\begin{itemize}
    \item \textbf{Sampler:} No-U-Turn Sampler (NUTS) \cite{hoffman2014nuts}
    \item \textbf{Chains:} 4 independent chains
    \item \textbf{Samples:} 2,000 draws per chain after 2,000 tuning steps
    \item \textbf{Convergence:} Gelman-Rubin $\hat{R} < 1.01$ for all parameters
    \item \textbf{Target acceptance:} 0.97 (high for stable geometry)
\end{itemize}

The posterior distribution $p(\theta_i | \mathbf{z})$ provides both a point estimate (posterior mean) and credible intervals (95\% HDI). For interpretability, we transform to a 0-100 scale using linear rescaling:

\begin{equation}
\text{Score}_{i} = 50 + 10 \cdot \mathbb{E}[\theta_i | \mathbf{z}]
\end{equation}

This preserves the z-score interpretation (mean=50, sd=10) while providing an intuitive range.

We apply this methodology to create four distinct composite scores, each targeting a specific prompt intent:

\begin{itemize}
    \item \texttt{crs\_100} - \textbf{Composite Reasoning Score:} Measures logical reasoning and problem-solving capabilities (0-100).
    \item \texttt{ccs\_100} - \textbf{Composite Coding Score:} Evaluates code generation and debugging proficiency (0-100).
    \item \texttt{cfs\_100} - \textbf{Composite Factual Score:} Assesses knowledge retrieval and factual accuracy (0-100).
    \item \texttt{css\_100} - \textbf{Composite Summarization Score:} Gauges ability to condense and synthesize information (0-100).
\end{itemize}

\subsection{Benchmark Weighting: Data-Driven vs. Manual}

Unlike weighted z-score approaches that require manual weight specification, BLF \textit{learns} benchmark importance from data via the loadings $\lambda_b$. Benchmarks with higher loadings contribute more to the composite score. Table \ref{tab:loadings} shows learned loadings for CCS:

\begin{table}[h]
\centering
\caption{Learned Benchmark Loadings for Composite Coding Score (CCS)}
\label{tab:loadings}
\begin{tabular}{lcccc}
\toprule
\textbf{Benchmark} & \textbf{$\lambda_b$} & \textbf{95\% HDI} & \textbf{$\sigma_b$} & \textbf{Coverage} \\
\midrule
LiveCodeBench     & 0.96 & [0.92, 1.01] & 0.29 & 90\% \\
HumanEval         & 0.88 & [0.83, 0.93] & 0.47 & 85\% \\
SciCode           & 0.91 & [0.86, 0.96] & 0.41 & 78\% \\
Arena Coding Rank$^*$ & 0.84 & [0.78, 0.90] & 0.54 & 82\% \\
Intelligence Index$^\dagger$ & 0.65 & [0.60, 0.70] & 0.76 & 99\% \\
\bottomrule
\end{tabular}
\begin{tablenotes}
\item[*] Inverted (lower rank = higher quality)
\item[$^\dagger$] Auxiliary benchmark for imputation
\end{tablenotes}
\end{table}

The high loadings ($\lambda > 0.8$) for primary coding benchmarks indicate strong agreement, while the auxiliary benchmark has lower loading, appropriately down-weighting it. The noise terms $\sigma_b$ reveal measurement reliability: HumanEval has higher noise (pass@1 is stochastic), while LiveCodeBench is more deterministic.

\subsection{Handling Missing Data: A Case Study}

Consider three models with different benchmark coverage:

\begin{itemize}
    \item \textbf{Model A:} All 4 primary benchmarks + auxiliary (complete data)
    \item \textbf{Model B:} 2 primary benchmarks + auxiliary (partial data)
    \item \textbf{Model C:} Only auxiliary benchmark (extreme missingness)
\end{itemize}

Figure \ref{fig:missing_data} shows how posterior uncertainty (HDI width) increases gracefully with missingness. Model C has wider credible intervals but still receives a meaningful score estimate, enabling inclusion in the ranking. Listwise deletion would discard Models B and C entirely.

\subsection{Model Validation}

We validate BLF against three baselines on held-out correlation with human preference data (Chatbot Arena ELO \cite{zheng2023lmsys}):

\begin{table}[h]
\centering
\caption{Correlation with Chatbot Arena ELO (Coding Category)}
\label{tab:validation}
\begin{tabular}{lcccc}
\toprule
\textbf{Method} & \textbf{Spearman $\rho$} & \textbf{Coverage} & \textbf{N Models} \\
\midrule
BLF (proposed)            & \textbf{0.89}$^{***}$ & 95\% & 247 \\
Weighted Z-Score          & 0.84$^{***}$ & 68\% & 177 \\
Arithmetic Mean           & 0.76$^{***}$ & 68\% & 177 \\
Best Single (LiveCodeBench) & 0.82$^{***}$ & 90\% & 234 \\
\bottomrule
\end{tabular}
\begin{tablenotes}
\item[***] $p < 0.001$
\end{tablenotes}
\end{table}

BLF achieves the highest correlation while maintaining near-complete coverage (95\% vs. 68\% for weighted z-score). The 5\% excluded models lack even auxiliary benchmarks.

\subsection{Convergence Diagnostics}

Figure \ref{fig:convergence} presents MCMC diagnostics for the CCS model:

\begin{itemize}
    \item \textbf{Trace plots} show good mixing across 4 chains with no drift
    \item \textbf{$\hat{R}$ statistics} are all $< 1.01$, indicating convergence
    \item \textbf{Effective sample size (ESS)} $> 4000$ for all parameters
    \item \textbf{Posterior predictive checks} show excellent agreement between observed and predicted z-scores (Figure \ref{fig:ppc})
\end{itemize}

These diagnostics confirm that the model has converged and is not misspecified.

\subsection{Computational Considerations}

\begin{itemize}
    \item \textbf{Inference time:} 3-5 minutes for 250 models $\times$ 5 benchmarks on a single CPU
    \item \textbf{Scaling:} $O(N \cdot M \cdot S)$ where $N$=models, $M$=benchmarks, $S$=MCMC samples
    \item \textbf{Caching:} Scores are computed offline and cached; online routing uses precomputed posteriors
    \item \textbf{Updates:} When new benchmark data arrives, only affected models need re-scoring
\end{itemize}

\subsection{Discussion}

The BLF model provides several advantages over naive aggregation:

\begin{enumerate}
    \item \textbf{Principled missing data handling:} No ad-hoc imputation or deletion required
    \item \textbf{Data-driven weighting:} Benchmark importance is learned, not manually specified
    \item \textbf{Uncertainty quantification:} Credible intervals inform routing decisions (e.g., prefer high-certainty models for risk-averse applications)
    \item \textbf{Robustness:} Outliers are naturally downweighted through Bayesian shrinkage
    \item \textbf{Interpretability:} Loadings reveal which benchmarks are most informative
\end{enumerate}

The primary limitation is computational cost: BLF requires MCMC sampling, whereas weighted z-scores are instantaneous. However, this is a one-time offline cost; online routing is fast. We also assume benchmarks are commensurable after standardization, which may not hold if they measure fundamentally different constructs (e.g., code generation vs. code review).

\subsection{Related Work}

BLF draws on classical factor analysis \cite{spearman1904general} and Item Response Theory (IRT) \cite{embretson2000item}, adapted to the LLM benchmarking context. It can be viewed as a continuous Bayesian IRT model. Similar latent variable approaches have been used in meta-learning \cite{hospedales2021metalearning} and multi-task learning \cite{caruana1997multitask}, but to our knowledge, this is the first application to LLM model selection with missing benchmark data.

The use of auxiliary benchmarks for imputation is inspired by multiple imputation \cite{rubin1987multiple} and Bayesian hierarchical models \cite{gelman2006data}. Our approach differs from standard multiple imputation by jointly modeling the latent factor and observed scores, avoiding separate imputation and analysis steps.
